\begin{figure}[t]
	\centering
	\begin{tikzpicture}[scale=0.95]
		% Define colors
		\colorlet{garnet}{Garnet}
		\colorlet{atlantic}{Atlantic}
		\colorlet{congaree}{Congaree}
		\definecolor{LightGrey}{RGB}{235,235,235}

		% ====================
		% LEFT PANEL: Automotive Radar
		% ====================
		\begin{scope}[shift={(0,0)}]
			% Panel background
			\draw[fill=LightGrey, draw=none] (-3.5,-3.5) rectangle (3.5,3.5);

			% Car outline (simplified rectangle)
			\draw[Garnet, line width=2pt] (-0.8,-0.5) rectangle (0.8,1.2);
			\draw[Garnet, line width=1.5pt] (-0.8,-0.5) circle (0.3);
			\draw[Garnet, line width=1.5pt] (0.8,-0.5) circle (0.3);

			% Radar cone (full FOV, quasi-simultaneous)
			\draw[Garnet, line width=2pt, fill=Garnet, fill opacity=0.15] (0,1.2) -- (2.5,3) -- (2.5,-3) -- cycle;
			\draw[Garnet, line width=2pt, fill=Garnet, fill opacity=0.15] (0,1.2) -- (-2.5,3) -- (-2.5,-3) -- cycle;

			% Radar antenna point
			\fill[Garnet] (0,1.2) circle (0.15);
			\draw[Garnet, line width=1.5pt] (0,1.2) -- (0,2);

			% Detection points at t₀
			\fill[Atlantic] (1.2,1.5) circle (0.12);
			\fill[Atlantic] (1.5,0) circle (0.12);
			\fill[Atlantic] (1.0,-1.8) circle (0.12);
			\fill[Atlantic] (-1.3,1.2) circle (0.12);
			\fill[Atlantic] (-1.6,-0.5) circle (0.12);
			\fill[Atlantic] (-0.8,-2) circle (0.12);

			% All detections labeled with same timestamp
			\node[Congaree, font=\small\bfseries] at (2.8,2.2) {$t_0$};
			\node[Congaree, font=\small\bfseries] at (2.8,0.5) {$t_0$};
			\node[Congaree, font=\small\bfseries] at (2.5,-1.2) {$t_0$};
			\node[Congaree, font=\small\bfseries] at (-2.8,1.8) {$t_0$};
			\node[Congaree, font=\small\bfseries] at (-2.8,-0.2) {$t_0$};
			\node[Congaree, font=\small\bfseries] at (-2.3,-2) {$t_0$};

			% Panel title
			\node[font=\bfseries\Large, Garnet] at (0,3.8) {Automotive Radar};
			\node[font=\bfseries, Atlantic] at (0,-3.8) {(Electronic Steering)};

			% Caption below panel
			\node[font=\small, text width=6cm, align=center] at (0,-4.5) {All targets measured at\\same time $t_0$};
		\end{scope}

		% ====================
		% RIGHT PANEL: Marine Radar
		% ====================
		\begin{scope}[shift={(8,0)}]
			% Panel background
			\draw[fill=LightGrey, draw=none] (-3.5,-3.5) rectangle (3.5,3.5);

			% Ship outline (simplified)
			\draw[Garnet, line width=2pt] (-1,-0.3) rectangle (1,1.2);
			\draw[Garnet, line width=1.5pt] (-0.3,0.4) -- (0.3,0.4);
			\draw[Garnet, line width=1.5pt] (-0.3,0.8) -- (0.3,0.8);

			% Rotating antenna mast
			\fill[Garnet] (0,1.2) circle (0.18);
			\draw[Garnet, line width=2pt] (0,1.2) -- (0,2.2);

			% Rotating radar antenna beam indicator
			\draw[Garnet, line width=2.5pt] (0,1.2) -- (2.2,1.8);

			% Pie slices showing different sweep times
			% Sector 1: t₀ (12 o'clock region)
			\draw[Garnet, line width=1.5pt, fill=Garnet, fill opacity=0.25] (0,1.2) -- (0,3) -- (1.3,2.6) -- cycle;
			\node[Congaree, font=\tiny\bfseries] at (0.5,2.5) {$t_0$};

			% Sector 2: t₀+Δ (3 o'clock region)
			\draw[Garnet, line width=1.5pt, fill=Garnet, fill opacity=0.20] (0,1.2) -- (1.3,2.6) -- (2.5,0.8) -- cycle;
			\node[Congaree, font=\tiny\bfseries] at (1.5,1.8) {$t_0+\Delta$};

			% Sector 3: t₀+2Δ (6 o'clock region)
			\draw[Garnet, line width=1.5pt, fill=Garnet, fill opacity=0.15] (0,1.2) -- (2.5,0.8) -- (1.5,-2) -- cycle;
			\node[Congaree, font=\tiny\bfseries] at (1.5,0) {$t_0+2\Delta$};

			% Sector 4: t₀+3Δ (9 o'clock region)
			\draw[Garnet, line width=1.5pt, fill=Garnet, fill opacity=0.10] (0,1.2) -- (1.5,-2) -- (-1.5,-2) -- cycle;
			\node[Congaree, font=\tiny\bfseries] at (0,-1.5) {$t_0+3\Delta$};

			% Sector 5: t₀+4Δ (back to 12 o'clock)
			\draw[Garnet, line width=1.5pt, fill=Garnet, fill opacity=0.05] (0,1.2) -- (-1.5,-2) -- (-1.3,2.6) -- cycle;
			\node[Congaree, font=\tiny\bfseries] at (-0.8,-0.5) {$t_0+4\Delta$};

			% Closing sector back to start
			\draw[Garnet, line width=1.5pt, fill=Garnet, fill opacity=0.00] (0,1.2) -- (-1.3,2.6) -- (0,3) -- cycle;

			% Camera frame indicator (at left, capturing at t₀)
			\draw[Atlantic, line width=2pt, dashed] (-3,-2.5) rectangle (-1.5,-0.5);
			\node[Atlantic, font=\small\bfseries] at (-2.25,-3.1) {Camera $t_0$};

			% Arrow showing mismatch
			\draw[Atlantic, line width=1.5pt, ->] (-1.3,-1.5) -- (-0.2,0.8);
			\node[Atlantic, font=\tiny] at (-1.5,-1.2) {Only};
			\node[Atlantic, font=\tiny] at (-1.5,-1.5) {this sector};

			% Panel title
			\node[font=\bfseries\Large, Garnet] at (0,3.8) {Marine Radar};
			\node[font=\bfseries, Atlantic] at (0,-3.8) {(Mechanical Rotation)};

			% Caption below panel
			\node[font=\small, text width=6cm, align=center] at (0,-4.5) {Different azimuths measured\\at different times};

			% Rotation period annotation
			\node[font=\tiny, Congaree, text width=5cm, align=center] at (0,-5.3) {1--3 sec full rotation\\$\rightarrow$ bearing-dependent timing error};
		\end{scope}

		% ====================
		% Main caption annotation (between panels)
		% ====================
		\node[font=\small, text width=2.5cm, align=center, Garnet] at (4,-6.5) {Temporal\\Mismatch};
		\draw[Garnet, line width=1.5pt, <->] (2.5,-6.2) -- (5.5,-6.2);

	\end{tikzpicture}
	\caption{Temporal synchronization challenge for mechanically rotating marine radar. Unlike electronically steered automotive radar that illuminates the entire FOV quasi-simultaneously, a rotating marine radar antenna measures different azimuths at different times within each sweep, creating bearing-dependent timing errors when associating with camera frames.}
	\label{fig:rotating_antenna_sync}
\end{figure}
